\chapter{Reference: text elements}
\label{ch:refman-text}

\begin{aside}
\code{text} for runtime strings; \code{t<"...">} for literals
known at compile time. Both produce \code{Element}.
\end{aside}

\section{Constructors}

\begin{itemize}[leftmargin=*]
  \item \code{text(std::string)} — owns the bytes.
  \item \code{text(std::string\_view)} — borrows; caller
    must outlive the frame.
  \item \code{text(int)}, \code{text(double)} — formatted as
    decimal.
  \item \code{t<"literal">} — zero-allocation form for
    compile-time constants.
\end{itemize}

\section{Working example}

\begin{lstlisting}[language=C++]
#include <maya/maya.hpp>

using namespace maya;
using namespace maya::dsl;

struct TextDemo {
    struct Model { int n = 0; std::string label = "live"; };
    struct Bump {}; struct Quit {};
    using Msg = std::variant<Bump, Quit>;

    static Model init() { return {}; }

    static auto update(Model m, Msg msg)
        -> std::pair<Model, Cmd<Msg>>
    {
        return std::visit(overload{
            [&](Bump) { m.n++; return std::pair{m, Cmd<Msg>{}}; },
            [](Quit)  { return std::pair{Model{}, Cmd<Msg>::quit()}; },
        }, msg);
    }

    static Element view(const Model& m) {
        return v(
            t<"compile-time literal">,
            text("runtime owned string"),
            text(m.label),                    // dynamic
            text(m.n),                        // formatted int
            text(std::format("{:.2f}", 3.14)) // your own format
        ) | pad<1> | border_<Round>;
    }

    static auto subscribe(const Model&) -> Sub<Msg> {
        return key_map<Msg>({
            {'q', Quit{}},
            {' ', Bump{}},
        });
    }
};

int main() { run<TextDemo>({.title = "text"}); }
\end{lstlisting}

\section{Notes}

\begin{itemize}[leftmargin=*]
  \item \code{t<"...">} is preferred whenever the string is
    a literal: it skips the allocation and the copy.
  \item For format strings, build the \code{std::string}
    first; pass the result to \code{text}.
  \item Unicode is supported; the layout engine treats wide
    glyphs as two cells.
\end{itemize}

\section{Recap}

\begin{recap}
\begin{itemize}[leftmargin=*]
  \item \code{t<>} for literals, \code{text} for runtime.
  \item Numeric overloads exist for convenience.
  \item Use \code{std::format} for anything fancy.
\end{itemize}
\end{recap}

\section*{Workshop}

\begin{enumerate}[leftmargin=*]
  \item Replace one \code{text(...)} with \code{t<"...">}
    where the value is a literal.
  \item Show the current time using
    \code{std::chrono::system\_clock}.
  \item Render the value of a \code{std::vector<int>}'s size
    via \code{text(int)}.
\end{enumerate}
